\documentclass[10pt]{article}
\usepackage[letterpaper,margin=.75in]{geometry}
\usepackage[T1]{fontenc}
\usepackage{lmodern,amsmath,amssymb,booktabs,graphicx,pdfpages,url,float}
\usepackage[hidelinks]{hyperref}
\setlength{\parindent}{0pt}
\setlength{\parskip}{6pt}
\providecommand{\reportbuilddatetime}{unknown}
\InputIfFileExists{report-build-info.tex}{}{}
\input{report-numbers.tex}
\title{Two-dimensional flattening of randomly sampled saddle graphs}
\author{\small Source: \path{~/current_projects/grip/papers/grip-software-paper/}\protect\\
\small\path{reproducibility/experiments/random-saddle/report.tex}}
\date{Build \reportbuilddatetime}
\begin{document}
\maketitle

\section*{Question and main finding}
Can symmetric nearest-neighbor graphs of randomly sampled saddle observations
be drawn in two dimensions while preserving their edge and fixed-path lengths?
This standalone study replaces the special orthogonal grid with independently
sampled observations. Nothing has been inserted into the manuscript.

The study completed \samplesTotal{} samples and \graphsTotal{} distinct weighted
graphs. A numerically robust four-vertex obstruction to exact two-dimensional
edge-length preservation was detected in \obstructionsTotal{} graphs. Nevertheless,
some two-dimensional candidates attain very small aggregate fixed-path errors.
These findings are compatible: exact zero loss would preserve every edge, whereas
a small average path error can coexist with larger local edge errors.

\section*{Sampling and graph construction}
The observations lie on $z=0.8(x^2-y^2)$ with $(x,y)\in[-1,1]^2$, without noise.
Sampling is uniform in surface area, not in the parameter square. Rejection
sampling uses a density proportional to $\sqrt{1+2.56(x^2+y^2)}$ in that square.
There are 100 independent samples at each of $n=250$ and $n=500$.

The \texttt{dgraphs} symmetric-kNN constructor uses exact Euclidean neighbors in
the original three-dimensional observations. An edge is retained when either
endpoint selects the other. No pruning or spanning-tree repair is applied.
Every integer $k$ is checked from 1 upward until the graph is connected; this
defines $k_{\rm conn}$. Each sample is then evaluated at $k_{\rm conn}$,
$k_{\rm conn}+2$, and $2k_{\rm conn}$. Edge lengths are divided by their median.
Connectedness is the selection rule, not a guarantee of surface-geodesic accuracy.

\section*{Coordinates and error criteria}
One input shortest path is retained for every unordered vertex pair and kept
fixed across configurations. All fitted configurations have exactly two
coordinate columns. The original three-dimensional observations give a known
zero-edge-loss and zero-path-loss reference.

For observed lengths $y_p$ and input lengths $g_p$, the scale is
$a=\sum_p y_pg_p/\sum_p g_p^2$, and the relative error is
\[
 R(y,g)=\left\{\frac{\sum_p(y_p-a g_p)^2}{\sum_p(a g_p)^2}\right\}^{1/2}.
\]
Edge and fixed-path summaries fit separate scales. The path criterion uses
embedded sums of edge lengths, not new shortest paths in the drawing. MDS
Stress-1 instead uses endpoint chords and the configuration denominator
$\sum_p y_p^2$. The plots report these dimensionless statistics as percentages.

\section*{Search budget and geometric obstruction}
Metric MDS, PCA of the original observations, and two seeded weighted-GRIP
configurations provide four starts. Each receives five-stage edge-KK refinement,
with at most 50 iterations per stage. Those four results and unrefined MDS and
PCA start six direct optimizations of squared profiled path relative RMSE, each
with at most 200 L-BFGS-B iterations. The lowest-path-error candidate is selected
from all 14 configurations. PCA uses the observations; this is a search for
low-distortion realizations, not a graph-only algorithm competition.

The best achieved error is an upper bound on the true optimum, not proof of
optimality. Independently, each four-vertex clique is checked using the Gram
matrix of its six squared edge lengths. A third/first eigenvalue ratio above
$10^{-8}$ is a robust numerical indication that those lengths require three
dimensions. The strongest witness is saved for every graph. This is not an
interval-arithmetic certificate. All six edges of each selected witness were
independently verified to be retained one-edge paths, so the obstruction also
excludes exact zero all-pairs fixed-path loss.

\section*{Repeated-sample results}
Table 1 reports the fixed-budget best-path candidate. Brackets give across-sample
interquartile ranges; they are not confidence intervals. Bootstrap 95\% intervals
for each median are in the accompanying CSV and figures. Edge errors belong to
the same selected configurations, with separately fitted scales.

\input{report-summary-table.tex}

At both sample sizes, the median achieved path error increases as the neighborhood
rule adds more neighbors. At the connectivity threshold the two sample sizes
have similar median path errors, whereas the denser rules have lower median path
errors at $n=500$. Edge errors of the selected path candidates are larger at
$n=500$ for all three rules. Thus increasing the number of observations does not
uniformly reduce all measures of distortion under this fixed optimization budget.

The neighborhood-rule comparison changes both the graph and its retained path
family. It is therefore a sensitivity analysis, not a comparison against one
unchanging metric. The distributions in Figures 1--3 distinguish the three
criteria and retain the sample-size and neighborhood-rule factors.

\section*{Optimization status and additional budget}
All attempted sample runs completed; none was excluded. The main study produced
\edgeFits{} fixed-budget edge-KK fits and \pathFits{} direct path fits.
Of the latter, \pathConverged{} satisfied the solver's numerical stopping rule,
\pathLimited{} reached the iteration limit, and \pathOther{} returned another
status. An iteration-limited result is retained as a finite candidate but is not
described as converged.

The winning candidate of a median-selected sample in each of the six groups
received up to 2,000 additional iterations. Table 2 separates this diagnostic
from the unchanged repeated-study results. A stopping criterion is not a global
optimality certificate, and these six selected cases do not estimate a new
population distribution.

\input{report-audit-table.tex}

\section*{Checks, figure selection, and interpretation}
Analytic gradients passed finite-difference checks. A planar control had zero
path error and no obstruction; a tetrahedron produced the expected obstruction.
The sampler was checked against numerical surface-area integration. All
\graphsTotal{} graphs were checked against independently recomputed shortest-path
distances and independent R summation of the saved paths. The maximum path-score
discrepancy in that post-run check was \pathCheckError{}.
In \nearTieGraphs{} graphs, near-tie handling differed from \texttt{igraph}; the
maximum scaled relative distance difference was \nearTieRelative{}, within the
package's $\sqrt{\text{machine epsilon}}$ tolerance. Substituting the stricter
distances changed the winner's path relative RMSE by at most \strictScoreChange{}.
One edge-endpoint pair used an indirect near-tied route; it did not belong to any
selected obstruction witness.

Figure 4 selects one sample per size by proximity to the median fixed-budget
best path error at $k_{\rm conn}$, breaking ties by replicate number. Selection
does not use appearance. Its two-dimensional layouts are displayed as aligned,
tilted planes; the scores use their original coordinates, not the projection.

The experiment establishes graph-level obstructions and achievable errors. It
does not measure continuous-surface geodesic error, prove recovery of curvature,
or guarantee that minimally connected graphs adequately represent the surface.
Most importantly, near-zero aggregate path error should not be interpreted as
near-zero local metric distortion. Exact zero and merely small loss carry
different implications here.

\includepdf[pages=-,fitpaper=true]{path-error-distributions.pdf}
\includepdf[pages=-,fitpaper=true]{edge-error-distributions.pdf}
\includepdf[pages=-,fitpaper=true]{chord-stress-distributions.pdf}
\includepdf[pages=-,fitpaper=true]{representative-layouts.pdf}
\end{document}
